% Приложение В — Ключевые фрагменты кода экспериментального фреймворка
\chapter{КЛЮЧЕВЫЕ ФРАГМЕНТЫ КОДА ЭКСПЕРИМЕНТАЛЬНОГО ФРЕЙМВОРКА}

В данном приложении приведены ключевые фрагменты исходного кода экспериментального фреймворка, реализованного на языке Python. Полный код опубликован в открытом доступе: \url{https://github.com/WpythonW/cognitive-biases-llm}. Для каждого фрагмента указан файл источника и номер строки в репозитории.

% -------------------------------------------------------
\section{Метрика IOU через метод Монте-Карло}
\label{sec:appC_iou}
% -------------------------------------------------------

Метрика Intersection-over-Union (IOU) используется в эксперименте~3 (задача 2-4-6) для оценки экстенсиональной эквивалентности гипотезы модели $H_t$ и истинного правила $R$. Оценка производится на выборке из $200\,000$ случайных троек $(x, y, z) \in [1, 1000]^3$.

Ниже приведены вспомогательные функции генерации выборки и безопасного вычисления лямбда-выражений (файл \texttt{wason-2-4-6/run\_experiments.py}):

\begin{lstlisting}
MIN_TEST_VALUE = 1
MAX_TEST_VALUE = 1000

def safe_eval(func, x, y, z):
    try:
        return bool(func(x, y, z))
    except Exception:
        return False

def generate_samples(n=50_000):
    return np.random.randint(
        MIN_TEST_VALUE, MAX_TEST_VALUE + 1, (n, 3)
    )
\end{lstlisting}

Основная функция \texttt{compare\_sets} вычисляет IOU, а также долю покрытия гипотезы в правиле ($H \subseteq R$) и правила в гипотезе ($R \subseteq H$):

\begin{lstlisting}
def compare_sets(R, H, n=200_000):
    samples = generate_samples(n)
    r_vec = np.array(
        [safe_eval(R, *s) for s in samples], dtype=bool
    )
    h_vec = np.array(
        [safe_eval(H, *s) for s in samples], dtype=bool
    )

    both = int((r_vec & h_vec).sum())
    h_in_r = both / max(int(h_vec.sum()), 1)
    r_in_h = both / max(int(r_vec.sum()), 1)
    dist = float(
        np.sqrt((1 - h_in_r)**2 + (1 - r_in_h)**2)
    )
    iou = compute_iou(r_vec, h_vec)
    return (
        round(h_in_r, 4), round(r_in_h, 4),
        round(dist, 4),   round(iou, 4),
    )

def compute_iou(r_vec, h_vec):
    intersection = int((r_vec & h_vec).sum())
    union = int((r_vec | h_vec).sum())
    return intersection / max(union, 1)
\end{lstlisting}

Испытание считается успешным при $\text{IOU}(H_t, R) > 0{,}95$.

% -------------------------------------------------------
\section{Метрика Severity of Bias}
\label{sec:appC_bias}
% -------------------------------------------------------

Метрика Severity of Bias (SoB) введена для преодоления ограничения бинарной классификации ошибок конъюнкции. Она использует самооценку уверенности модели для измерения \emph{силы} предубеждения. Файл \texttt{linda-problem/run\_experiments.py}.

Извлечение выбора и уверенности из ответа модели осуществляется с помощью регулярных выражений:

\begin{lstlisting}
CHOICE_PATTERN = re.compile(
    r"^\s*(?:ВЫБОР|CHOICE)\s*:\s*[\[\(]?\s*"
    r"([abAB])\s*[\]\)]?\s*$",
    re.MULTILINE,
)
CONFIDENCE_PATTERN = re.compile(
    r"^\s*(?:УВЕРЕННОСТЬ|CONFIDENCE)\s*:\s*"
    r"\[?\s*(100|[0-9]{1,2})\s*\]?\s*$",
    re.MULTILINE,
)

def parse_model_answer(content):
    choice_match = CHOICE_PATTERN.search(content)
    confidence_match = CONFIDENCE_PATTERN.search(content)
    if choice_match is None or confidence_match is None:
        raise ValueError(
            f"Could not parse: {content!r}"
        )
    choice = choice_match.group(1).lower()
    confidence = int(confidence_match.group(1))
    if choice not in {"a", "b"}:
        raise ValueError(f"Unexpected choice: {choice!r}")
    if not 0 <= confidence <= 100:
        raise ValueError(
            f"Confidence out of range: {confidence}"
        )
    return ParsedAnswer(
        choice=choice, confidence=confidence
    )
\end{lstlisting}

Вычисление метрики SoB: при ошибочном выборе конъюнкции (b) значение пропорционально уверенности; при нормативно верном выборе (a) --- обратно пропорционально:

\begin{lstlisting}
def compute_bias_score(choice, confidence):
    if choice == "b":
        return confidence / 100.0
    if choice == "a":
        return (100 - confidence) / 100.0
    raise ValueError(f"Unsupported choice: {choice!r}")
\end{lstlisting}

Таким образом, $\text{SoB} \in [0, 1]$, где 0 соответствует нормативно верному выбору с максимальной уверенностью, а 1 --- ошибочному выбору с абсолютной уверенностью. Разность $\Delta\text{bias} = \overline{\text{SoB}}_{\text{corr}} - \overline{\text{SoB}}_{\text{uncorr}}$ измеряет вклад нарративной согласованности в силу предубеждения.

% -------------------------------------------------------
\section{Протокол парсинга и ремонта ответов}
\label{sec:appC_parsing}
% -------------------------------------------------------

Для каждого эксперимента реализован протокол парсинга ответов модели с механизмом автоматического ремонта при нарушении формата.

\subsection{Задача Линды}

Парсер извлекает структурированные поля (CHOICE, CONFIDENCE) из свободного текста ответа модели с помощью регулярных выражений, допускающих вариации форматирования (скобки, квадратные скобки, регистр). Соответствующий код приведён в разделе~В.2.

\subsection{Задача выбора карточек Уэйсона}

Парсер извлекает список карточек из текста ответа, обрабатывая формат \texttt{ANSWER:} для условия Chain-of-Thought (файл \texttt{wason-selection/run\_experiments.py}):

\begin{lstlisting}
def parse_response(response_text, json_mapping):
    normalized_text = response_text.strip()
    if "\n" in normalized_text:
        non_empty = [
            l.strip() for l in
            normalized_text.splitlines()
            if l.strip()
        ]
        if non_empty:
            normalized_text = non_empty[-1]
    if normalized_text.upper().startswith("ANSWER:"):
        normalized_text = (
            normalized_text.split(":", 1)[1].strip()
        )
    cards = [
        c.strip()
        for c in normalized_text.split(",")
        if c.strip()
    ]
    logic_set = frozenset(
        marker for marker in (
            json_mapping.get(c) for c in cards
        ) if marker is not None
    )
    return cards, logic_set
\end{lstlisting}

Функция \texttt{json\_mapping} переводит текстовые метки карточек в логические маркеры ($P$, $\lnot P$, $Q$, $\lnot Q$), что позволяет оценивать правильность выбора в абстрактном пространстве логических ролей.

\subsection{Задача 2-4-6: парсинг и протокол ремонта}

Парсер извлекает гипотезу ($H$) и тестовую тройку ($T$) из fenced-блока Python (файл \texttt{wason-2-4-6/run\_experiments.py}):

\begin{lstlisting}
def parse_response(text):
    stripped = text.strip()
    if stripped.startswith("```"):
        lines_raw = stripped.splitlines()
        if (len(lines_raw) >= 3
            and lines_raw[0].startswith("```")
            and lines_raw[-1].strip() == "```"):
            stripped = "\n".join(lines_raw[1:-1]).strip()

    lines = [l.strip() for l in stripped.splitlines()
             if l.strip()]
    hypothesis = next(
        (l[2:].strip() for l in lines
         if l.startswith("H:")), None
    )
    test_str = next(
        (l[2:].strip() for l in lines
         if l.startswith("T:")), None
    )
    if not hypothesis or not test_str:
        raise ValueError(f"Parse failed: {text}")

    test_obj = eval(test_str)
    if not isinstance(test_obj, list) or len(test_obj) != 3:
        raise ValueError(f"Invalid triple: {test_obj}")

    test = [int(v) for v in test_obj]
    if any(v < 1 or v > 1000 for v in test):
        raise ValueError(
            f"Values out of [1..1000]: {test}"
        )
    return hypothesis, test
\end{lstlisting}

При некорректном ответе модели отправляется ремонтный промпт, указывающий на конкретную ошибку:

\begin{lstlisting}
def build_repair_prompt(error):
    return (
        "Your previous answer was invalid.\n"
        f"Error: {error}\n"
        "Re-answer ONLY with a fenced "
        "```python``` block containing:\n"
        "H: lambda x,y,z: <bool expr>\n"
        "T: [x, y, z]\n"
        "Use only integers from 1 to 1000."
    )
\end{lstlisting}

Протокол допускает до \texttt{repair\_attempts} попыток ремонта. Если все попытки исчерпаны, испытание помечается как неуспешное.

% -------------------------------------------------------
\section{Статистический анализ}
\label{sec:appC_stats}
% -------------------------------------------------------

Статистический протокол включает точный тест МакНемара для согласованных пар, парный t-тест, критерий Уилкоксона и процедуру Бенджамини-Хохберга для контроля частоты ложных открытий (FDR). Файл \texttt{linda-problem/compute\_metrics.py}.

\subsection{Точный тест МакНемара}

Тест МакНемара применяется для сравнения частоты бинарных ошибок в парных условиях (correlated vs.\ uncorrelated). Используется точный биномиальный вариант:

\begin{lstlisting}
def mcnemar_exact(n12, n21):
    discordant = n12 + n21
    if discordant == 0:
        exact_pvalue = 1.0
    else:
        exact_pvalue = float(
            binomtest(
                min(n12, n21), discordant, 0.5
            ).pvalue
        )
    statistic = (
        0.0 if discordant == 0
        else ((abs(n12 - n21) - 1)**2) / discordant
    )
    return {
        "n12": float(n12), "n21": float(n21),
        "discordant": float(discordant),
        "mcnemar_statistic_cc": float(statistic),
        "mcnemar_exact_pvalue": exact_pvalue,
    }
\end{lstlisting}

\subsection{Процедура Бенджамини-Хохберга}

Процедура контролирует ожидаемую долю ложных открытий при множественном тестировании:

\begin{lstlisting}
def apply_bh(pvalues):
    if not pvalues:
        return []
    indexed = sorted(enumerate(pvalues),
                     key=lambda item: item[1])
    adjusted = [0.0] * len(pvalues)
    running_min = 1.0
    total = len(pvalues)
    for rank, (orig_idx, pval) in enumerate(
        reversed(indexed), start=1
    ):
        adj = min(1.0, pval * total / (total - rank + 1))
        running_min = min(running_min, adj)
        adjusted[orig_idx] = running_min
    return adjusted
\end{lstlisting}

\subsection{Парные тесты для метрики SoB}

Для сравнения непрерывных значений Severity of Bias между условиями применяются парный t-тест и критерий Уилкоксона:

\begin{lstlisting}
def paired_bias_tests(pair_df):
    corr = pair_df["bias_score_correlated"].astype(float)
    uncorr = pair_df["bias_score_uncorrelated"].astype(float)

    t_stat, t_pvalue = ttest_rel(corr, uncorr)
    nonzero = (corr - uncorr) != 0
    if nonzero.any():
        w_stat, w_pvalue = wilcoxon(
            corr[nonzero], uncorr[nonzero],
            zero_method="wilcox"
        )
    else:
        w_stat, w_pvalue = 0.0, 1.0

    return {
        "mean_bias_correlated": float(corr.mean()),
        "mean_bias_uncorrelated": float(uncorr.mean()),
        "delta_bias": float(
            corr.mean() - uncorr.mean()
        ),
        "paired_t_statistic": float(t_stat),
        "paired_t_pvalue": float(t_pvalue),
        "wilcoxon_statistic": float(w_stat),
        "wilcoxon_pvalue": float(w_pvalue),
    }
\end{lstlisting}

% -------------------------------------------------------
\section{Конвейер генерации датасета}
\label{sec:appC_dataset}
% -------------------------------------------------------

Генерация датасета задачи Линды включает два этапа: создание биографической виньетки и генерацию меток (профессия T, хобби F). Файл \texttt{linda-problem/generate\_dataset.py}.

Функция \texttt{compose\_row} собирает итоговую строку датасета из виньетки и меток:

\begin{lstlisting}
def compose_row(vignette, labels):
    t = clean_phrase(labels.t)
    f = clean_phrase(labels.f)
    f_uncorr = clean_phrase(labels.f_uncorrelated)
    return {
        "vignette": vignette,
        "T": t,
        "F": f,
        "F_uncorrelated": f_uncorr,
        "T&F1": f"{t} and {f}",
        "T&F2": f"{t} and {f_uncorr}",
    }
\end{lstlisting}

Асинхронная генерация виньетки по заданной профессии с автоматическими повторами при ошибке:

\begin{lstlisting}
async def generate_persona(client, seed_job, semaphore):
    async with semaphore:
        last_error = None
        for _ in range(MAX_ROW_ATTEMPTS):
            try:
                parsed = await request_structured(
                    client,
                    model=MODEL,
                    messages=build_persona_messages(
                        seed_job,
                        None if last_error is None
                        else str(last_error)
                    ),
                    response_format=PersonaOutput,
                    schema_name="linda_persona_output",
                    temperature=TEMPERATURE,
                )
                return normalize_vignette(parsed.vignette)
            except Exception as exc:
                last_error = exc
        raise ValueError(
            f"Failed: seed_job={seed_job!r}: "
            f"{last_error}"
        )
\end{lstlisting}

Генерация меток с валидацией уникальности хобби:

\begin{lstlisting}
async def generate_labels(client, vignette, semaphore,
                          forbidden_hobbies):
    async with semaphore:
        last_error = None
        for _ in range(MAX_ROW_ATTEMPTS):
            try:
                parsed = await request_structured(
                    client,
                    model=MODEL,
                    messages=build_label_messages(
                        vignette, forbidden_hobbies,
                        None if last_error is None
                        else str(last_error)
                    ),
                    response_format=LabelOutput,
                    schema_name="linda_label_output",
                    temperature=TEMPERATURE,
                )
                if (clean_phrase(parsed.f).lower()
                    == clean_phrase(
                        parsed.f_uncorrelated).lower()):
                    raise ValueError(
                        "f and f_uncorrelated identical"
                    )
                return parsed
            except Exception as exc:
                last_error = exc
        raise ValueError(
            f"Failed to generate labels: {last_error}"
        )
\end{lstlisting}

Каждая виньетка проходит верификацию: проверка уникальности коррелированного и некоррелированного хобби, дедупликация по семантическому пересечению с ранее сгенерированными строками.
